\documentclass[notes=show]{beamer}
%%%%%%%%%%%%%%%%%%%%%%%%%%%%%%%%%%%%%%%%%%%%%%%%%%%%%%%%%%%%%%%%%%%%%%%%%%%%%%%%%%%%%%%%%%%%%%%%%%%%%%%%%%%%%%%%%%%%%%%%%%%%%%%%%%%%%%%%%%%%%%%%%%%%%%%%%%%%%%%%%%%%%%%%%%%%%%%%%%%%%%%%%%%%%%%%%%%%%%%%%%%%%%%%%%%%%%%%%%%%%%%%%%%%%%%%%%%%%%%%%%%%%%%%%%%%
\AtBeginSubsection[
  {\frame<beamer>{\frametitle{Outline}   
    \tableofcontents[currentsection,currentsubsection]}}%
]%
{
  \frame<beamer>{ 
    \frametitle{Outline}   
    \tableofcontents[currentsection,currentsubsection]}
}
\usepackage{ulem}
\usepackage{pdflscape}
\setbeamertemplate{caption}{\insertcaption}
\usepackage{color,xcolor,ucs}
\usepackage{beamerthemesplit}
\usepackage{graphicx}
\usepackage{amsmath}
\usepackage{adjustbox}
\usepackage{pdfpages} 
\definecolor{mintgreen}{rgb}{0.6, 1.0, 0.6}
\makeatletter
\@addtoreset{subfigure}{framenumber}% subfigure counter resets every frame
\makeatother
\newcommand*{\Scale}[2][4]{\scalebox{#1}{$#2$}}%
\newcommand*{\Resize}[2]{\resizebox{#1}{!}{$#2$}}%
\usepackage{amsfonts}
\usepackage{booktabs}
\usepackage{color,soul}
\usepackage{pdfpages}
\setboolean{@twoside}{false}
\usepackage{amsmath}
\usepackage{pdflscape} 
\usepackage{eurosym}
\usepackage{amstext} % for \text
\DeclareRobustCommand{\officialeuro}{%
  \ifmmode\expandafter\text\fi
  {\fontencoding{U}\fontfamily{eurosym}\selectfont e}}
\usepackage[caption = false]{subfig}
\usepackage{appendixnumberbeamer} 
\usepackage{graphicx}
\usepackage{mathpazo}
\usepackage{hyperref}
\usepackage{multimedia}
\usepackage{graphicx}
\usepackage{multirow}
\usepackage{subfig}
\usepackage{verbatim}
\usepackage{booktabs, multicol, multirow}
\setcounter{MaxMatrixCols}{10}
\input{Macros.tex}
\AtBeginSection[]
{
  \begin{frame}[noframenumbering]
    \frametitle{Outline for section \thesection}
    \tableofcontents[currentsection]
  \end{frame}
}


\usepackage{xcolor,soul}
\renewcommand<>{\hl}[1]{\only#2{\beameroriginal{\hl}}{#1}}

% https://tex.stackexchange.com/questions/41683/why-is-it-that-coloring-in-soul-in-beamer-is-not-visible
\makeatletter
\newcommand\SoulColor{%
  \let\set@color\beamerorig@set@color
  \let\reset@color\beamerorig@reset@color}
\makeatother
\SoulColor
\AtBeginSection{\frame{\sectionpage}}

\newsavebox\mybox
\newenvironment{aquote}[1]
  {\savebox\mybox{#1}\begin{quote}\openautoquote\hspace*{-.7ex}}
  {\unskip\closeautoquote\vspace*{1mm}\signed{\usebox\mybox}\end{quote}}

\newtheorem{proposition}[theorem]{Proposition}
\newtheorem{Assumption}[theorem]{Assumption}
\newenvironment{stepenumerate}{\begin{enumerate}[<+->]}{\end{enumerate}}
\newenvironment{stepitemize}{\begin{itemize}[<+->]}{\end{itemize} }
\newenvironment{stepenumeratewithalert}{\begin{enumerate}[<+-| alert@+>]}{\end{enumerate}}
\newenvironment{stepitemizewithalert}{\begin{itemize}[<+-| alert@+>]}{\end{itemize} }
\usetheme{Madrid}
\makeatletter
\setbeamertemplate{footline}
{
  \leavevmode%
  \hbox{%
  \begin{beamercolorbox}[wd=.333333\paperwidth,ht=2.25ex,dp=1ex,center]{author in head/foot}%
    \usebeamerfont{author in head/foot}\insertshortauthor~~\beamer@ifempty{\insertshortinstitute}{}{(\insertshortinstitute)}
  \end{beamercolorbox}%
  \begin{beamercolorbox}[wd=.333333\paperwidth,ht=2.25ex,dp=1ex,center]{title in head/foot}%
    \usebeamerfont{title in head/foot}\insertshorttitle
  \end{beamercolorbox}%
  \begin{beamercolorbox}[wd=.333333\paperwidth,ht=2.25ex,dp=1ex,right]{date in head/foot}%
    \usebeamerfont{date in head/foot}\insertshortdate{}\hspace*{2em}
%    \insertframenumber{} / \inserttotalframenumber\hspace*{2ex} % DELETED
  \end{beamercolorbox}}%
  \vskip0pt%
}
\usepackage[style=british]{csquotes}

\def\signed #1{{\leavevmode\unskip\nobreak\hfil\penalty50\hskip1em
  \hbox{}\nobreak\hfill #1%
  \parfillskip=0pt \finalhyphendemerits=0 \endgraf}}

\makeatother
\usepackage{pdfpages}
\begin{document}
	\title[]{\large Search and Unemployment		\vspace{0.5cm}
	}
	\author[Raffaele Saggio - Econ 560]{Raffaele Saggio}
	\institute[]{UBC}
	\date{\today \\}
	\maketitle

%%%%%%%%%%%%%%%%%%%%%%%%%	
%%%%%%%%%%%%%%%%%%%%%%%%%
%%%%%%%%%%%%%%%%%%%%%%%%%
%%%%%%%%%%%%%%%%%%%%%%%%%
\setbeamertemplate{footline}[frame number]{}

%gets rid of bottom navigation symbols
\setbeamertemplate{navigation symbols}{}

%gets rid of footer
%will override 'frame number' instruction above
%comment out to revert to previous/default definitions
%%%%%%%%%%%%%%%%%%%%%%%%%%%%%
\begin{frame}
\begin{itemize}
\item What is unemployment?
\medskip
\item Concept arose in the 1930s from the US statistics agency in the aftermath of the Great depression.
\medskip
\end{itemize}
\end{frame}
\begin{frame}
\frametitle{Measurament}
\centering
\includegraphics[width=1\textwidth]{aux/card_U}
\end{frame}
%%%%%%%%%%%%%%%%%%%%%%%%%%%%%
\begin{frame}
\frametitle{Getting/Losing a Job is a BIG DEAL}
\centering
\includegraphics[width=1\textwidth]{aux/manning}
\end{frame}
%%%%%%%%%%%%%%%%%%%%%%%%%%%%%
\setbeamertemplate{footline}{}
\begin{frame}
\frametitle{Why there is unemployment?}
\begin{itemize}
\item Assume an intertemporal labor supply function of the form (Lucas and Rapping, 1969). 
\begin{equation}
log h_{it} = \alpha_{i} + x_{it}'\beta + \eta log w_{it} + \epsilon_{it}
\end{equation}
\pause
\begin{equation}
log h^{*}_{it} = \alpha_{i} + x_{it}'\beta + \eta log \underbrace{w^{*}_{it}}_{\text{expected wages}} + r_{it}
\end{equation}
\item This leads to one potential notion of unemployment
\be
u_{it}=\lambda(log h^{*}_{it} - log h_{it}) = \lambda\eta  (log {w^{*}_{it}} - log w_{it}) + \lambda(\epsilon_{it}-r_{it})
\ee
\pause
\item When wages are lower than expected, retract labor supply and that goes into leisure.
\medskip
\item (Some) Macro guys believe that $\eta=4$ $\rightarrow$ small reduction in wages leads to take a vacation in Greece for a month. 
\medskip
\item No space for the concept of ``unemployed" defined as someone actively searching for a job (and is unable to find it!).
\end{itemize}
\end{frame}
%%%
%%
\begin{frame}
\frametitle{Unemployment $\approx$ Leisure in Neo-classical Labor Supply model}
\begin{figure}[htb]
\centering
\includegraphics[width=0.8\textwidth]{aux/hurst}
\end{figure}
\end{frame}
%%
\begin{frame}
\frametitle{Measurament ahead of Theory}
\begin{figure}[htb]
\centering
\includegraphics[width=0.8\textwidth]{aux/card_U2.jpg}
\end{figure}
\end{frame}
%%
\begin{frame}
\frametitle{Search}
\begin{itemize}
\item Search models are an alternative (and possibly more reasonable) way to explain why there is unemployment.
\medskip
\item Plan:
\medskip
\begin{itemize}
\item Brief discussion of search models and concept of the reservation wage.
\medskip
\item The role of Unemployment insurance.
\medskip
\item The long-term consequences of becoming unemployed.
\end{itemize}
\end{itemize}
\end{frame}
%%%
\section{Search and General Facts About Unemployment}
\begin{frame}[plain]{McCall (1970)}
\begin{itemize}
\item Individuals enter the labor market as unemployed and receive some benefit while unemployed given by $b$. They live forever and are risk neutral with discount factor $\beta = \dfrac{1}{1+r}<1$. 
\medskip
\item Each period they search and receive a job offer that can be accepted/rejected. Search has (monetary) cost $\nu$.
\medskip
\item Once job is accepted it lasts forever. There is no on the job search (and utility is linear in money). 
\medskip
\item Q: what is the p.d.v of working? 
\end{itemize}
\end{frame}

\begin{frame}[shrink=10]
\frametitle{Value Functions}
\begin{itemize}
\item Let $U$ denote the optimal utility of a worker who is however currently unemployed. 
\medskip
\item When will a worker accept a job?
\pause
\item Wages are drawn from \underline{exogenous} offer distribution with cdf $F\left(.\right)$ and pdf $f(.)$.
\begin{equation}
\underbrace{U}_{\text{Value of being unemployed}} = b-\nu + \beta \int_{0}^{\infty}\max\{U,\dfrac{w}{1-\beta}\}f(w)d(w)
\end{equation}
\item Now define $w^{*}$ has the level of wage such that $\dfrac{w^{*}}{1-\beta}=U$
\medskip
\item We call $w^{*}$ the reservation wage: lowest level of wage that will make the worker exit unemployment state.
\pause
\medskip
\item Useful expression that implicity defines the reservation wage:
\begin{equation}
w^{*}= b-\nu + \dfrac{1}{r}\E[w-w^{*}\vert w\geq w^{*}]Pr(w\geq w^{*}) 
\end{equation}
\end{itemize}
\end{frame}
\begin{frame}{Unemployment}
\begin{itemize}
    \item All workers here have the same reservation wage.
    \item Therefore, the probability to leave the unemployment state is given by 
\begin{equation}
Pr(w\geq w^{*})=1-F(w^{*})
\end{equation}
\item Now the policy question: what happens when we increase unemployment benefits?
\begin{equation}
    \dfrac{\partial w^{*}}{\partial b} = \dfrac{1}{1+\dfrac{1}{r}(1-F(w^{*}))}>0
\end{equation}


\pause
\item Hence, generous unemployment benefits should make searchers picky, raise accepted wages, and lengthen
unemployment durations.
\end{itemize}
\end{frame}
\begin{frame}{Comments}
\begin{itemize}
    \item Constant reservation wage.
    \medskip
    \item Hard to get data on offer distribution.
    \medskip
    \item Partial equilibrium model (why?).
\end{itemize}
\end{frame}

\begin{frame}[plain]{Krueger and Mueller (2016)}

\begin{itemize}
\item How important are reservation wages, really?
\medskip
\item Survey roughly 6,000 unemployed workers in New Jersey over the period
October 2009 through March 2010. 
\medskip
\item Data available \href{http://opr.princeton.edu/archive/njui/.\%20}{online}
but low (10\%) response rate
\medskip
\item Study how reservation wages should vary with unemployment duration
using non-stationary search model ala Mortensen (1977).
\item Find that actual behavior is quite different. 
\end{itemize}
\end{frame}

\begin{frame}[plain]{Model}

\begin{itemize}
\item {\scriptsize{}The value function of an unemployed worker with $t$
months of benefits left (there are $T$ months of benefits that the worker can have) is:
\[
U\left(t\right)=u\left(b\left(t\right)\right)+\beta\max_{R}\left\{ U\left(t-1\right)+\alpha\int_{R}\left(W\left(x,m=0\right)-U\left(t-1\right)\right)dF\left(x\right)\right\} 
\]
where $R$ is the reservation wage, and $W\left(w,m=0\right)$ is
the value of starting a job paying wage $w$. }{\scriptsize\par}
\item {\scriptsize{}To qualify for UI benefits, workers must stay with their
employer for $\bar{m}=6$ months. The value in month $m<\bar{m}$
of a job paying wage $w$ is:
\begin{eqnarray*}
W\left(w,m\right) & = & u\left(w\right)+\beta\left[\left(1-\delta\right)W\left(w,m+1\right)+\delta U\left(0\right)\right]\\
 &  & +\beta\left[\alpha_{e}\left(1-\delta\right)\int_{w}\left(W\left(x,m+1\right)-W\left(w,m+1\right)\right)dF\left(x\right)\right]
\end{eqnarray*}
where $\alpha_{e}$ is the probability of finding another job while
employed. }{\scriptsize\par}
\item {\scriptsize{}In months, $m\geq\bar{m}$, the worker has qualified
for benefits and the problem becomes stationary with:
\begin{eqnarray*}
W\left(w,\bar{m}\right) & = & u\left(w\right)+\beta\left[\left(1-\delta\right)W\left(w,\bar{m}\right)+\delta U\left(T\right)\right]\\
 &  & +\beta\alpha_{e}\left(1-\delta\right)\int_{w}\left(W\left(x,\bar{m}\right)-W\left(w,\bar{m}\right)\right)dF\left(x\right)
\end{eqnarray*}
}{\scriptsize\par}
\end{itemize}
\end{frame}

\begin{frame}[plain]{Model}

\begin{itemize}
\item {\footnotesize{}The reservation wage function $R\left(t\right)$ solves
the following equation:
\[
W\left(R\left(t\right),m=0\right)=U\left(t-1\right)
\]
}{\footnotesize\par}
\item {\footnotesize{}No closed form, solve on computer}{\footnotesize\par}
\end{itemize}
\centering
\includegraphics[scale=0.3]{aux/KM_fig1.png}
\end{frame}

\begin{frame}[plain]{Doesn't seem to work!}

\begin{center}
\centering
\includegraphics[scale=0.4]{aux/KM_fig2.png}
\par\end{center}

\end{frame}

\begin{frame}[plain]{Look at the level!}

\begin{center}
\centering
\includegraphics[scale=0.4]{aux/KM_fig2A.png}
\par\end{center}

\end{frame}

\begin{frame}[plain]{Focusing on the sample of those that received job offers...}

\begin{center}
\includegraphics[scale=0.4]{aux/KM_fig3.png}
\par\end{center}

\end{frame}
\begin{frame}[plain]{Focusing on the sample of those that received job offers...}

\begin{center}
\includegraphics[scale=0.4]{aux/KM_fig3_bis}
\par\end{center}
e
\end{frame}

\begin{frame}[plain]{Diverse motives for rejecting offers above reservation}

\begin{center}
\includegraphics[scale=0.4]{aux/KM_fig4.png}
\par\end{center}

\end{frame}

\begin{frame}[plain]{What is going on?}

\begin{center}
\includegraphics[scale=0.4]{aux/KM_fig5.png}
\par\end{center}

\end{frame}
\begin{frame}[plain]{Summing up}

\begin{itemize}
\item Limited impact of UI benefits on reservation wages.
\pause
\bigskip
\item But why? Any suggestions?
\par\end{itemize}

\end{frame}
\section{Policy Debate over UI}
\begin{frame}
\includegraphics[scale=0.2]{aux/shortage}
\end{frame}
\begin{frame}
\includegraphics[scale=0.4]{aux/UI_FT}
\end{frame}


\begin{frame}[plain]{Schmeider, Von Wachter, and Bender (2016)}

\begin{itemize}
\item In standard search models, lengthening/increasing benefit duration should raise
reservation wages (thus higher unemployment duration) but also lead to higher re-employment wages
\medskip
\item An earlier experimental literature (Meyer, 1996; Robins and Spiegelman,
2001) finds that hurrying the search process up doesn't do much to
the re-employment wage

\begin{itemize}
\item Some of this could just be imprecision
\end{itemize}
\item Recent studies use quasi-experiments with massive datasets to look
for a precise effect.
\item In Germany, maximum benefit duration increases discontinuously with
age, leads to a natural RDD to test for impacts on re-employment wage
\end{itemize}
\end{frame}
\begin{frame}{Context}
\begin{itemize}
    \item If involuntary separated from the employer at age 42 or younger: 12 months of maximum coverage of UI ($\approx$ 63\% replacement rate if without children, 68\% with children)
    \medskip
    \item If involuntary separated from the employer at age 43 or 44: 18 months of UI.
    \medskip
    \item At age 45 further increase in maximum duration to 22 months. 
\end{itemize}
    
\end{frame}

\begin{frame}[plain]{Strong first stage effect on durations}

\begin{center}
\includegraphics[scale=0.4]{aux/SvWB_fig2.png}
\par\end{center}

\end{frame}

\begin{frame}[plain]{Extra time = (a) just leisure or (b) search for better jobs?}

\begin{center}
\includegraphics[scale=0.35]{aux/SvWB_fig2.png}
\par\end{center}

\end{frame}

\begin{frame}[plain]{Small \emph{negative} impact on re-employment wage!}

\begin{center}
\includegraphics[scale=0.4]{aux/SvWB_fig3.png}
\par\end{center}

\end{frame}
\begin{frame}[plain]{Making sense of the evidence...}

\begin{itemize}
\item Negative impact on re-employment wages hard to reconcile with standard search models.
\medskip
\item Note that these models assume that nothing truly happens to the offer distribution if you stay 12 months unemployed as opposed to only 1 month.
\medskip
\item This could unrealistic: skills might decay the more time you spent unemployed. Even if skills don't actually decay, more time spent unemployed can send a negative signal to employers.
\medskip
\item This skill decay / duration dependence argument can reconcile the evidence in Schmeider et al. (2016). But do we have any evidence on this?
\end{itemize}

\end{frame}

\begin{frame}[plain]{Making sense of the evidence...}

\begin{itemize}
\item Negative impact on re-employment wages hard to reconcile with standard search models.
\medskip
\item Note that these models assume that nothing truly happens to the offer distribution if you stay 12 months unemployed as opposed to only 1 month.
\medskip
\item This could unrealistic: skills might decay the more time you spent unemployed. Even if skills don't actually decay, more time spent unemployed can send a negative signal to employers.
\medskip
\item This skill decay / duration dependence argument can reconcile the evidence in Schmeider et al. (2016). But do we have any evidence on this?
\end{itemize}
\end{frame}
\begin{frame}[plain]{}

\includegraphics[scale=0.25]{aux/noto2.jpg}
\end{frame}
\begin{frame}[plain]{}

\includegraphics[scale=0.3]{aux/noto3.jpg}
\end{frame}

\begin{frame}[plain]{Nekoei and Weber (2017)}

\begin{itemize}
\item Examine an age-based discontinuity in maximum benefit duration in
Austria
\item They find that longer durations \emph{increase} the re-employment
wage
\item Reconciliation: they argue that duration dependence sometimes outweighs
reservation wage effects
\item Provide evidence for this argument via meta-analysis of relationship
between duration and wage impacts
\end{itemize}
\end{frame}

\begin{frame}[plain]{Small increase in the wage}

\includegraphics[scale=0.4]{aux/Nekoei_Weber_fig3}
\end{frame}

\begin{frame}[plain]{A meta-analysis}

\begin{center}
\includegraphics[scale=0.3]{aux/Nekoei_Weber_figAA}
\par\end{center}

\end{frame}

\begin{frame}[plain]{A meta-analysis}

\begin{center}
\includegraphics[scale=0.35]{aux/Nekoei_Weber_fig4}
\par\end{center}

\end{frame}

\begin{frame}[plain]{Thoughts}

\begin{itemize}
\item Across all studies of this subject, wage effects are very small (but also first stage)
\medskip
\begin{itemize}
\item NW: increase in 10 weeks of UI $\rightarrow$ 2 extra days job less (2\% increase) $\rightarrow$ 0.5 percent higher wage
\end{itemize}
\medskip
\item Nekoei and Weber write a directed search model:
\begin{enumerate}
\medskip
\item UI induces workers to search for better jobs.
\medskip
\item But you can also become more picky (search of a perfect job) $\rightarrow$ more unemployed $\rightarrow$ negative duration dependence. 
\end{enumerate}
\medskip
\item Heterogeneity in the matching function can generate  negative relationship between effect of wages and UI unemployment duration. 
\end{itemize}
\end{frame}
\begin{frame}
\frametitle{More recent Evidence}
\centering
\includegraphics[scale=0.25]{aux/stepner0}
\end{frame}
\begin{frame}

\frametitle{Treated and Control}
\centering
\includegraphics[scale=0.3]{aux/stepner1}
\end{frame}
\begin{frame}
\frametitle{Effects on Employment}
\centering
\includegraphics[scale=0.3]{aux/stepner2}
\end{frame}
\begin{frame}
\frametitle{An awesome figure}
\centering
\includegraphics[scale=0.3]{aux/stepner3}
\end{frame}
\begin{frame}[plain]{Krueger and Mueller (2010)}

\begin{itemize}
\item What does it mean to be unemployed? How can we measure search intensity?
\item Official definition: ``people are classified as unemployed if they
do not have a job, have actively looked for work in the prior 4 weeks,
and are currently available for work.''
\item Krueger and Mueller study to see what nominally unemployed individuals
are doing with their time.
\item Data from ATUS from 2003-2007. 
\end{itemize}
\end{frame}

\begin{frame}[plain]{Job search does not seem time intensive}

\begin{center}
\includegraphics[scale=0.4]{aux/KM10_fig1.png}
\par\end{center}

\end{frame}

\begin{frame}[plain]{Unemployed search more than the employed}

\begin{center}
\includegraphics[scale=0.4]{aux/KM10_fig2}
\par\end{center}

\end{frame}

\begin{frame}[plain]{Benefits strongly correlated with time use}

\begin{center}
\includegraphics[scale=0.4]{aux/KM10_fig3.png}
\par\end{center}

\end{frame}

\begin{frame}[plain]{Procrastination?}

\begin{center}
\includegraphics[scale=0.6]{aux/KM10_fig4.png}
\par\end{center}
\end{frame}

\begin{frame}[plain]{2020 Version of this plot: Marinescu and Skandalis (2020)}

\begin{center}
\includegraphics[scale=0.18]{aux/marinescu_skandalis}
\par\end{center}
\end{frame}

\begin{frame}[plain]{2020 Version of this plot: DellaVigna, Heining, Schmieder and Trenkle (2020)}
\begin{center}
\includegraphics[scale=0.18]{aux/johannes}
\par\end{center}
\end{frame}

\begin{frame}[plain]{Unemployment and Well-Being}

\begin{center}
\includegraphics[scale=0.5]{aux/kruegerWELLBEING.png}
\par\end{center}
\end{frame}

\section{The Consequences of losing a job}
\begin{frame}
\frametitle{Now let's look at things from a different angle}
\begin{itemize}
\item What are the consequences of a job loss \underline{for an individual perspective?}
\bigskip
\item Neo-classical theory: I will always get paid my marginal product of labor.
\bigskip
\item Things are of course much more complicated (lost of firm specific skills, stigma, rents, etc).
\end{itemize}
\end{frame}
\begin{frame}{Jacobson, LaLonde and Sullivan (1993)}
\begin{itemize}
\item {\small{}Jacobson, LaLonde and Sullivan (1993) look at the effects
of job displacement on workers' earnings trajectories\medskip{}
}{\small \par}
\item {\small{}JLS use matched administrative records on workers and firms
in Pennsylvania to look at what happens when high-tenure (6 or more
years) workers lose their jobs\medskip{}
}{\small \par}
\item {\small{}Worker data: Unemployment insurance records\medskip{}
}{\small \par}
\item {\small{}Firm data: ES202 forms (filed by all firms covered by UI)\medskip{}
}{\small \par}
\item {\small{}To focus on involuntary job separations, construct a ``mass-layoff''
sample of separations that occured at firms during times when employment
dropped by 30\% or more}{\small \par}
\end{itemize}
\end{frame}

\begin{frame}{Jacobson, LaLonde and Sullivan (1993)}

\begin{itemize}
\item {\small{}JLS equation:}{\small \par}
\end{itemize}
\begin{center}
{\small{}$y_{it}=\alpha_{i}+\gamma_{t}+x_{it}'\beta+{\displaystyle \sum_{k\geq-m}^{M}D_{it}^{k}\delta_{k}+\epsilon_{it}}$}
\par\end{center}{\small \par}

\end{frame}

\begin{frame}{Jacobson, LaLonde and Sullivan (1993)}

\begin{itemize}
\item {\small{}JLS equation:}{\small \par}
\end{itemize}
\begin{center}
{\small{}$y_{it}=\alpha_{i}+\gamma_{t}+x_{it}'\beta+{\displaystyle \sum_{k\geq-m}^{M}D_{it}^{k}\delta_{k}+\epsilon_{it}}$.}
\par\end{center}{\small \par}
\begin{itemize}
\item {Let $t^{*}_{i}$ denote the year of displacement of individual $i$.
}{\small \par}

\item {\small{}$D_{it}^{k}$: one if individual $i$ was displaced exactly
$k$ periods before $t$\medskip{}
}{\small \par}
\be
D_{it}^{k} = \mathbf{1}\{t=t^{*}_{i} + k\}
\ee
\medskip
\item {\small{}Since some $k$'s are negative, this specification also looks
at earnings dynamics before separation $\rightarrow$ pre-trends\medskip{}
}{\small \par}
\item {\small{}Worker and time effects $\implies$ this is a differences
in differences estimator\medskip{}
}{\small \par}
\pause
\item Who is the control group here?
\medskip
\item {\small{}$\delta_{k}$ coefficients compare changes in earnings for
workers whose displacement status changed to changes for workers who
were not displaced at the same time}{\small \par}
\end{itemize}
\end{frame}
\begin{frame}
\begin{figure}[htb]
\centering
\includegraphics[width=0.95\textwidth]{aux/jls_f1}
\end{figure}
\end{frame}
\begin{frame}
\begin{figure}[htb]
\centering
\includegraphics[width=0.95\textwidth]{aux/jls_f2}
\end{figure}
\end{frame}
\begin{frame}
\begin{figure}[htb]
\centering
\includegraphics[width=0.95\textwidth]{aux/jls_f3}
\end{figure}
\end{frame}
\begin{frame}{Jacobson, LaLonde and Sullivan (1993)}

\begin{itemize}
\item {\small{}JLS findings:\medskip{}
}

\begin{itemize}
\item {\small{}Relative earnings for workers who will experience mass layoffs
start to decline in advance, then drop sharply at displacement\medskip{}
}{\small \par}
\item {\small{}Quick partial recovery, but stabilize around 25\% below initial
level\medskip{}
}{\small \par}
\item {\small{}Younger workers experience larger losses, but faster recovery\medskip{}
}{\small \par}
\item {\small{}Larger losses in mining/construction/public utilities; perhaps
loss of high-rent union jobs?\medskip{}
}{\small \par}
\item {\small{}Long-term losses even for workers reemployed in same industry.
JLS interpret this as evidence of firm-specific skills...}{\small \par}
\end{itemize}
\end{itemize}
\end{frame}

\begin{frame}{Praise of JLS}

\begin{itemize}
\item Incredibly important paper for empirical labor economics
\begin{enumerate}
\bigskip
\item Empirical tool (Event study design).
\bigskip
\item Importance of the question.
\bigskip
\item Use of high quality administrative data.
\bigskip
\item Transparent design $\rightarrow$ replication in other countries and for different outcomes.
\end{enumerate}
\end{itemize}

\end{frame}

\begin{frame}[shrink=10]
\frametitle{Other Studies of Job Displacement}

\begin{itemize}
\item {\small{}Several recent papers by von Wachter and co-authors use JLS-style
methods to look at job displacement in other contexts\medskip{}
}

\begin{itemize}
\item von Wachter, Song and Manchester (2009): Use social security records
to study long-term effects of displacement in 1982 recession (more
states, 20 years of follow-up){\small{}\medskip{}
}{\small \par}
\item von Wachter and Handwerker (2009): Use California UI records matched
to CPS interviews. Allows examination of heterogeneity by education{\small{}\medskip{}
}{\small \par}
\item Sullivan and von Wachter (2009): Match JLS Pennsylvania data to Social
Security death records to look at mortality{\small{}\medskip{}
}{\small \par}
\item von Wachter, Bender and Schmeider (2009): Use German social security
data to look at 1982 recession in Germany. Includes detailed information
(e.g. days worked) on each employment spell, so allows an analysis
of wages
\end{itemize}
\end{itemize}
\end{frame}

\begin{frame}
\begin{figure}[htb]
\centering
\includegraphics[width=0.95\textwidth]{aux/till_1_f3}
\end{figure}
\end{frame}

\begin{frame}
\begin{figure}[htb]
\centering
\includegraphics[width=0.95\textwidth]{aux/till_2_f1}
\end{figure}
\end{frame}

\begin{frame}
\begin{figure}[htb]
\centering
\includegraphics[width=0.8\textwidth]{aux/till_3_f2}
\end{figure}
\end{frame}


\begin{frame}
\begin{figure}[htb]
\centering
\includegraphics[width=0.95\textwidth]{aux/till_4_f6}
\end{figure}
\end{frame}

\begin{frame}
\begin{figure}[htb]
\centering
\includegraphics[width=0.95\textwidth]{aux/daily_wages}
\end{figure}
\end{frame}

\begin{frame}
\frametitle{Evidence from Europe: Bertheau et al. (2023) }
\begin{figure}[htb]
\centering
\includegraphics[width=1\textwidth]{aux/jls_europe.jpg}
\end{figure}
\end{frame}
\begin{frame}
\frametitle{Evidence from Europe: Bertheau et al. (2023)}
\begin{figure}[htb]
\centering
\includegraphics[width=1\textwidth]{aux/jls_europe2.jpg}
\end{figure}
\end{frame}
\begin{frame}
\frametitle{Discussion: Why these enormous effects after losing a job?}
\begin{itemize}
\item Ideas? 
\pause
\medskip
\medskip
\item Jarosch JMP is a nice attempt: structural model that allows for heterogenous jobs in both security and wages. 
\medskip
\item Unemployed workers fell of the ladder: sorted into worst jobs + possible depreciation of skills + lost of bargaining chip.  
\pause
\includegraphics[width=0.6\textwidth]{aux/gregor}
\end{itemize}
\end{frame}
\begin{frame}
\frametitle{But labor market institutions are also important!}
\centering
\includegraphics[width=0.9\textwidth]{aux/jls_europe3}
\end{frame}

\begin{frame}{Discussion}

\begin{itemize}
\item {\small{}Displacement costs are real.}{\small \par}
\item {\small{}Need to re-think carefully about impact of labor demand shocks and/or policies.}{\small \par}
\item {\small{} What are the labor market consequences to interval/external shocks that might destroy jobs? }{\small \par}
\item {\small{} One example: China Shock. \medskip{}
}{\small \par}
\item {\small{} Another useful example: Green Policies.}
\end{itemize}
\end{frame}
%%
\begin{frame}{Walker (2013)}
  \begin{figure}[htb]
\centering
\includegraphics[width=0.95\textwidth]{aux/reed.pdf}
\end{figure}  
\end{frame}
%%
\begin{frame}{Effects of Child-Birth: Kleven-Landais-Sogaard, 2019}
\pause
\begin{figure}[htb]
\centering
\includegraphics[width=0.95\textwidth]{aux/kleven_child.jpg}
\end{figure}
\end{frame}
\begin{frame}{No sign of convergence 20 years after}
\begin{figure}[htb]
\centering
\includegraphics[width=0.95\textwidth]{aux/kleven_child2.jpg}
\end{figure}
\end{frame}
\begin{frame}{Culture Matters?}
\begin{figure}[htb]
\centering
\includegraphics[width=0.95\textwidth]{aux/kleven_child3.jpg}
\end{figure}
\end{frame}
\begin{frame}{Evidence from various countries}
\begin{figure}[htb]
\centering
\includegraphics[width=0.95\textwidth]{aux/kleven_child4.jpg}
\end{figure}
\end{frame}
\begin{frame}{Evidence from various countries}
\begin{figure}[htb]
\centering
\includegraphics[width=0.95\textwidth]{aux/kleven_child6.jpg}
\end{figure}
\end{frame}
\begin{frame}{Evidence from various countries}
\begin{figure}[htb]
\centering
\includegraphics[width=0.95\textwidth]{aux/kleven_child5.jpg}
\end{figure}
\end{frame}
\begin{frame}{Chilbirth of Women Entrepeneurs spillovers into the performance of the firm}
\begin{figure}[htb]
\centering
\includegraphics[width=0.65\textwidth]{aux/rutigliano.jpg}
\end{figure}
\end{frame}
\begin{frame}{Much more muted effects for men!}
\begin{figure}[htb]
\centering
\includegraphics[width=0.65\textwidth]{aux/rutigliano2.jpg}
\end{figure}
\end{frame}



%%
\begin{frame}
\centering
\LARGE{Appendix}
\end{frame}
\begin{frame}
\label{derivation}
\frametitle{A simple Search Model}
\begin{itemize}
\item Time is discrete. 
\item Unemployed individual who is infinitely lived, risk neutral, with discount factor: $\beta=\dfrac{1}{1+r}$
\item Each period receives a job offer that she can accept/reject.
\item If job is accepted, it lasts forever. There is no on the job search.
\item Individual employed at wage $w$ as flow utility  $w$.
\item Unemployed individual receives benefit $b$ and has flow utility $\nu$.
\item Wages are drawn from distribution $F(w)$ with pdf $f(w)$.
\end{itemize}
\end{frame}
%%
\begin{frame}[shrink=10]
\frametitle{Value Function / Bellman Equation}
\begin{itemize}
\item We can set up the bellman equation that defines the value of unemployment
\be
V = b - \nu + \beta \int_{0}^{\infty}max(V,\dfrac{w}{1-\beta})f(w)dw
\ee
\item This expression can be written as
\be
\begin{aligned}
V &   = b - \nu + \beta \int_{0}^{\infty}\left(V+max(0,\dfrac{w}{1-\beta}-V)\right)f(w)dw \\
V (1-\beta)  &   =   b - \nu + \dfrac{\beta}{1-\beta} \int_{0}^{\infty}\left(max(0,w-(1-\beta)V)\right)f(w)dw 
\end{aligned}
\ee

\item What is the reservation wage? $\dfrac{w^{*}}{1-\beta}=V$
\item Plugging this expression we then get:
\be
w^{*} = b-\nu + \dfrac{1}{r}\int_{0}^{\infty}(w-w^{*})f(w)dw 
\ee
\item which can be restated as:
\be
w^{*}-(b-\nu)=Pr(w\geq w^{*})\times \dfrac{\E[w-w^{*}\vert w>w^{*}]}{r}
\ee
which coincides with equation (\ref{key_Eq})
\end{itemize}
\end{frame}
\begin{frame}{Jarosch (2015)}

\begin{itemize}
\item {\small{}Jarosch (2015) approaches job displacement with a fresh theoretical
eye\medskip{}
}{\small \par}
\item {\small{}The long-term ``scarring'' effect of job displacement is
a major reason we're concerned with unemployment\medskip{}
}{\small \par}
\item {\small{}But most modern workhorse models used to study unemployment
(e.g. DMP search models) do not generate an unemployment scar\medskip{}
}{\small \par}
\item {\small{}Jarosch proposes a search model that can generate long-term
effects of job displacement\medskip{}
}{\small \par}
\item {\small{}Then applies the model to the German social security data
studied by von Wachter and co.}{\small \par}
\end{itemize}
\end{frame}

\begin{frame}{Jarosch (2015): Model}

\begin{itemize}
\item {\footnotesize{}Jarosch proposes a ``job ladder'' model with on-the-job
search \medskip{}
}{\footnotesize \par}
\item {\footnotesize{}Workers search for new jobs while employed or unemployed\medskip{}
}{\footnotesize \par}
\item {\footnotesize{}Jobs are heterogeneous on two dimensions:\medskip{}
}

\begin{itemize}
\item {\footnotesize{}Productivity}{\footnotesize \par}
\item {\footnotesize{}Job security (expected length of match)\medskip{}
}{\footnotesize \par}
\end{itemize}
\item {\footnotesize{}Workers ``climb'' the ladder, searching for better
and better jobs \textendash{} implies unemployment begets future unemployment
spells\medskip{}
}{\footnotesize \par}
\item {\footnotesize{}Model also allows employment history to alter workers'
bargaining position, and general human capital to evolve differentially
for employed and unemployed\medskip{}
}{\footnotesize \par}
\item {\footnotesize{}Three mechanisms here that can generate persistent
displacement effects: position in the job ladder, experience, bargaining
position}{\footnotesize \par}
\end{itemize}
\end{frame}

\begin{frame}{Setup}

\begin{itemize}
\item {\footnotesize{}Firm types: $\theta=[\theta_{y},\theta_{\delta}]$\medskip{}
}

\begin{itemize}
\item $\theta_{y}$ is productivity
\item $\theta_{\delta}$ is job distruction probability (1 - job security){\footnotesize{}\medskip{}
}{\footnotesize \par}
\end{itemize}
\item Workers are homogeneous \emph{ex ante}{\footnotesize{}\medskip{}
}{\footnotesize \par}
\item Unemployed workers meet firms with probability $\lambda_{0}$, employed
workers $\lambda_{1}$
\end{itemize}
\end{frame}

\begin{frame}{Wages and Value Functions}

\begin{itemize}
\item {\footnotesize{}Fixed wage contracts, re-bargained when either party
has a credible threat (sequential auction framework of Cahuc, Postel-Vinay and Robin, 2006)\medskip{}
}{\footnotesize \par}
\item {\footnotesize{}$w\left(\theta,\hat{\theta}\right)$: Worker's wage
when matched with type $\theta$, and $\hat{\theta}$ used as outside
option in last negotiation (negotiation benchmark)\medskip{}
}{\footnotesize \par}
\item {\footnotesize{}Notation for value functions:}

\begin{itemize}
\item {\footnotesize{}$W(\theta,\hat{\theta})$: employed worker's value}{\footnotesize \par}
\item {\footnotesize{}$U$: unemployed worker's value}{\footnotesize \par}
\item {\footnotesize{}$J$: value of a job}{\footnotesize \par}
\item {\footnotesize{}$S(\theta)$: joint surplus of a match between a worker
and firm $\theta$}{\footnotesize \par}
\end{itemize}
\end{itemize}
\end{frame}

\begin{frame}{Wages}

\begin{itemize}
\item {\footnotesize{}For bargaining parameter $\alpha\in[0,1]$, wage for
a newly employed worker at firm type $\theta$ will be negotiated
to satisfy}{\footnotesize \par}
\end{itemize}
\begin{center}
{\footnotesize{}$W(\theta,u)-U=\alpha S(\theta)$}
\par\end{center}{\footnotesize \par}
\begin{itemize}
\item {\footnotesize{}Suppose a worker at firm type $\theta_{1}$ is offered
a job at outside firm type $\theta_{2}$\medskip{}
}{\footnotesize \par}
\item {\footnotesize{}If $S(\theta_{2})>S(\theta_{1})$, the worker transfers
to $\theta_{2}$ and $\theta_{1}$ becomes the negotiation benchmark\medskip{}
}{\footnotesize \par}
\item {\footnotesize{}New wage $w(\theta_{2},\theta_{1})$ is negotiated
to satisfy}{\footnotesize \par}
\end{itemize}
\begin{center}
{\footnotesize{}$W(\theta_{2},\theta_{1})-U=S(\theta_{1})+\alpha\left(S(\theta_{2})-S(\theta_{1})\right)$}
\par\end{center}{\footnotesize \par}
\begin{itemize}
\item {\footnotesize{}Worker gains a fraction $\alpha$ of the extra surplus
generated by the better match}{\footnotesize \par}
\end{itemize}
\end{frame}

\begin{frame}{Wages}

\begin{itemize}
\item {\footnotesize{}If $S(\theta_{2})<S(\theta_{1})$ , the worker will
not transfer, but if $\theta_{2}$ is better than previous negotiation
benchmark may renegotiate wage at $\theta_{1}$ so that}{\footnotesize \par}
\end{itemize}
\begin{center}
{\footnotesize{}$W(\theta_{1},\theta_{2})-U=S(\theta_{2})+\alpha\left(S(\theta_{1})-S(\theta_{2})\right)$}
\par\end{center}{\footnotesize \par}
\begin{itemize}
\item {\footnotesize{}This suggests the following classification of firms:\medskip{}
}

\begin{itemize}
\item {\footnotesize{}$M_{1}(\theta)$: Firm types $x$ such that $S(x)>S(\theta)$}{\footnotesize \par}
\item {\footnotesize{}$M_{2}(\theta,\hat{\theta})$: Firm types $x$ such
that $S(\theta)>S(x)>S(\hat{\theta})$ \medskip{}
}{\footnotesize \par}
\end{itemize}
\item {\footnotesize{}A worker at $\theta$ with negotiation benchmark $\hat{\theta}$
will move if offered something in $M_{1}(\theta)$, stay but renegotiate
if offered something in $M_{2}(\theta,\hat{\theta})$, and ignore
if offered anything else}{\footnotesize \par}
\end{itemize}
\end{frame}

\begin{frame}{Negotiation Capital}

\begin{itemize}
\item {\footnotesize{}In the sequential auction framework, workers build
``negotiation capital'' through repeated sampling of outside offers\medskip{}
}{\footnotesize \par}
\item {\footnotesize{}Formally, negotiation capital is $W(\theta,\hat{\theta})-W(\theta,u)$\medskip{}
}{\footnotesize \par}
\item {\footnotesize{}This is the extra value a worker gets from $\theta$
because s/he is negotiating based on the benchmark $\hat{\theta}$
rather than unemployment\medskip{}
}{\footnotesize \par}
\item {\footnotesize{}This can generate large cross-sectional wage dispersion\medskip{}
}{\footnotesize \par}
\item {\footnotesize{}Also has the potential to generate sharp wage reductions
following an unemployment spell, as negotiation capital is destroyed
when a worker enters unemployment}{\footnotesize \par}
\end{itemize}
\end{frame}

\begin{frame}{Value Function for an Employed Worker}

\begin{center}
{\footnotesize{}$W(\theta,\hat{\theta})=w(\theta,\hat{\theta})+\beta\times\{\ (1-\theta_{\delta})$}
\par\end{center}{\footnotesize \par}

\begin{center}
{\footnotesize{}$\times[\ \lambda_{1}\left({\displaystyle \int_{x\in M_{1}(\theta)}W(x,\theta)dF(x)+\int_{x\in M_{2}(\theta,\hat{\theta)}}W(\theta,x)dF(x)}\right)$}
\par\end{center}{\footnotesize \par}

\begin{center}
{\footnotesize{}$+\left(1-\lambda_{1}{\displaystyle \int_{x\in M_{1}(\theta)\cup M_{2}(\theta,\hat{\theta})}dF(x)}\right)W(\theta,\hat{\theta})\ ]+\theta_{\delta}U\ \}$.}
\par\end{center}{\footnotesize \par}
\begin{itemize}
\item {\footnotesize{}Flow utility $w(\theta,\hat{\theta})$ , discount
rate $\beta$}{\footnotesize \par}
\item {\footnotesize{}With probability $\theta_{\delta}$, the match ends
and the worker enters unemployment}{\footnotesize \par}
\item {\footnotesize{}With probability $(1-\theta_{\delta})$, the match
doesn't end, and}

\begin{itemize}
\item {\footnotesize{}With probability $\lambda_{1}$, the worker meets
a new firm}

\begin{itemize}
\item {\footnotesize{}In $M_{1}(\theta)$: Switch and get $W(x,\theta)$}{\footnotesize \par}
\item {\footnotesize{}In $M_{2}(\theta,\hat{\theta}):$ Don't switch but
renegotiate for $W(\theta,x)$}{\footnotesize \par}
\item {\footnotesize{}In neither: Stick with $W(\theta,\hat{\theta})$}{\footnotesize \par}
\end{itemize}
\item {\footnotesize{}With probability $1-\lambda_{1}$, the worker does
not meet a firm and keeps $W(\theta,\hat{\theta})$}{\footnotesize \par}
\end{itemize}
\end{itemize}
\end{frame}

\begin{frame}{Value Function for an Unemployed Worker}

\begin{center}
{\footnotesize{}$U=z+\beta\left[\lambda_{0}{\displaystyle \int_{x\in M_{1}(u)}}W(x,u)dF(x)+\left(1-\lambda_{0}{\displaystyle \int_{x\in M_{1}(u)}dF(x)}\right)U\right]$ }
\par\end{center}{\footnotesize \par}
\begin{itemize}
\item {\footnotesize{}Flow utility is unemployment benefit $z$}{\footnotesize \par}
\item {\footnotesize{}With probability $\lambda_{0}$, the worker meets
a new firm}

\begin{itemize}
\item {\footnotesize{}In $M_{1}(u)$: Take the job and get $W(x,u)$}{\footnotesize \par}
\item {\footnotesize{}Not in $M_{1}(u)$: Reject job, keep $U$}{\footnotesize \par}
\end{itemize}
\item {\footnotesize{}With probability $1-\lambda_{0}$, the worker doesn't
meet a new firm, and keeps $U$}{\footnotesize \par}
\end{itemize}
\end{frame}

\begin{frame}{Value Function for a Firm}

\begin{center}
{\footnotesize{}$J(\theta,\hat{\theta})=\theta_{y}-w(\theta,\hat{\theta})+\beta(1-\theta_{\delta})\times[\ \lambda_{1}{\displaystyle \int_{x\in M_{2}(\theta,\hat{\theta})}J(\theta,x)dF(x)}$}
\par\end{center}{\footnotesize \par}

\begin{center}
{\footnotesize{}$+\left(1-\lambda_{1}{\displaystyle \int_{x\in<M_{1}(\theta)\cup M_{2}(\theta,\hat{\theta})}dF(x)}\right)J(\theta,\hat{\theta})\ ]$}
\par\end{center}{\footnotesize \par}
\begin{itemize}
\item {\footnotesize{}Flow utility for a firm of type $\theta$ employing
a worker with benchmark $\hat{\theta}$ is productivity $\theta_{y}$
minus wage $w(\theta,\hat{\theta})$}{\footnotesize \par}
\item {\footnotesize{}With probability $\theta_{\delta}$, job is destroyed
and has no continuation value}{\footnotesize \par}
\item {\footnotesize{}With probability $(1-\theta_{\delta})$, job continues,
and}

\begin{itemize}
\item {\footnotesize{}With probability $\lambda_{1}$, worker finds a new
job}

\begin{itemize}
\item {\footnotesize{}In $M_{1}(\theta)$: Job destroyed, no continuation
value}{\footnotesize \par}
\item {\footnotesize{}In $M_{2}(\theta,\hat{\theta})$: Renegotiate, get
$J(\theta,x)$}{\footnotesize \par}
\item {\footnotesize{}In neither: Stick with $J(\theta,\hat{\theta})$}{\footnotesize \par}
\end{itemize}
\item {\footnotesize{}With probability $1-\lambda_{1}$, worker does not
find a job and firm keeps $J(\theta,\hat{\theta})$}{\footnotesize \par}
\end{itemize}
\end{itemize}
\end{frame}

\begin{frame}{Match Surplus}

\begin{itemize}
\item {\footnotesize{}Match surplus for the worker/firm pair is}{\footnotesize \par}
\end{itemize}
\begin{center}
{\footnotesize{}$S(\theta,\hat{\theta})={\displaystyle \max\left\{ 0,W(\theta,\hat{\theta})-U+J(\theta,\hat{\theta})\right\} }$}
\par\end{center}{\footnotesize \par}
\begin{itemize}
\item {\footnotesize{}Plugging in value functions yields}{\footnotesize \par}
\end{itemize}
\begin{center}
{\footnotesize{}$S(\theta)=\theta_{y}-z+\beta\times\{\ (1-\theta_{\delta})\left(S(\theta)+\alpha\lambda_{1}\int_{x\in M_{1}(\theta)}\left(S(x)-S(\theta)\right)dF(x)\right)$}
\par\end{center}{\footnotesize \par}

\begin{center}
{\footnotesize{}$-\alpha\lambda_{0}{\displaystyle \int_{x\in M_{1}(u)}S(x)dF(x)\}}$}
\par\end{center}{\footnotesize \par}
\begin{itemize}
\item {\footnotesize{}Result: Surplus doesn't depend on $\hat{\theta}$,
so $S(\theta,\hat{\theta})=S(\theta)$}{\footnotesize \par}
\item {\footnotesize{}Makes sense \textendash{} negotiation benchmark just
determines split of the surplus}{\footnotesize \par}
\item {\footnotesize{}This equation can be solved numerically, given parameter
values}{\footnotesize \par}
\end{itemize}
\end{frame}

\begin{frame}{Comparative Statics}

\begin{itemize}
\item {\footnotesize{}Model gives rise to some intuitive comparative statics}

\begin{itemize}
\item {\footnotesize{}Wages increasing in $\theta_{\delta}$ (compensating
differential)}{\footnotesize \par}
\item {\footnotesize{}$\theta_{y}$ has an ambiguous effect on wages, depending
on $\alpha$}

\begin{itemize}
\item {\footnotesize{}When $\alpha=0$, moving workers get exactly their
old match's surplus. Higher productivity firms offer higher growth
potential and therefore lower wages}{\footnotesize \par}
\item {\footnotesize{}When $\alpha=1$, workers capture the entire surplus
and higher productivity firms pay higher wages}{\footnotesize \par}
\end{itemize}
\end{itemize}
\item {\footnotesize{}If $E\left[\theta_{\delta}|\theta_{y}\right]$ and
$E\left[\theta_{y}|\theta_{\delta}\right]$ are strictly increasing:}

\begin{itemize}
\item {\footnotesize{}Job security and productivity are increasing in employment
tenure and job tenure}{\footnotesize \par}
\item {\footnotesize{}Hazard to unemployment is decreasing in employment
tenure and job tenure}{\footnotesize \par}
\end{itemize}
\end{itemize}
\end{frame}

\begin{frame}{Efficiency}

\begin{itemize}
\item {\footnotesize{}Jarosch assumes $\lambda_{1}<\lambda_{0}$. This implies
the model's equilibrium is inefficient if $\alpha<1$}{\footnotesize \par}
\end{itemize}
\end{frame}

\begin{frame}{Efficiency}

\begin{itemize}
\item {\footnotesize{}Jarosch assumes $\lambda_{1}<\lambda_{0}$. This implies
the model's equilibrium is inefficient if $\alpha<1$\medskip{}
}{\footnotesize \par}
\item {\footnotesize{}Social planner places more value on unemployment than
does a matched worker/firm pair, because unemployed workers are available
to match with other firms, which is not internalized\medskip{}
}{\footnotesize \par}
\item {\footnotesize{}Workers therefore overvalue job security, relative
to the social optimum\medskip{}
}{\footnotesize \par}
\item {\footnotesize{}As a result, a positive unemployment benefit financed
by a lump sum tax on employed workers can improve welfare}{\footnotesize \par}
\end{itemize}
\end{frame}

\begin{frame}{Worker Skills}

\begin{itemize}
\item {\footnotesize{}Jarosch extends the model to add worker heterogeneity\medskip{}
}{\footnotesize \par}
\item {\footnotesize{}Workers have skill $s$, so productivity is $p(\theta_{y},s)$\medskip{}
}

\begin{itemize}
\item {\footnotesize{}Increases probabilistically for employed workers}{\footnotesize \par}
\item {\footnotesize{}Depreciates probabilisitically for unemployed workers\medskip{}
}{\footnotesize \par}
\end{itemize}
\item {\footnotesize{}Still no permanent }\emph{\footnotesize{}ex ante}{\footnotesize{}
differences in workers}{\footnotesize \par}
\end{itemize}
\end{frame}

\begin{frame}{Estimation}

\begin{itemize}
\item {\footnotesize{}Jarosch uses the model to quantify the importance
of the job ladder, negotiation capital and experience effects in generating
the displacement ``scar''\medskip{}
}{\footnotesize \par}
\item {\footnotesize{}Matches moments from German social security data,
1974-2010\medskip{}
}{\footnotesize \par}
\item {\footnotesize{}Method of simulated moments estimator fitting various
moments related to model parameters, e.g.:}

\begin{itemize}
\item {\footnotesize{}Unemployment rate ($\lambda_{0}$)}{\footnotesize \par}
\item {\footnotesize{}Job-to-job transition rate ($\lambda_{1})$}{\footnotesize \par}
\item {\footnotesize{}Unemployment hazard and slope w.r.t. tenure (distribution
of $\theta_{\delta}$)}{\footnotesize \par}
\item {\footnotesize{}Relationship between wage and length of unemployment
duration, w/worker FE (skill depreciation)}{\footnotesize \par}
\item {\footnotesize{}Relationship between wages and separation rate (correlation
of $\theta_{\delta}$ and $\theta_{y}$)}{\footnotesize \par}
\end{itemize}
\end{itemize}
\end{frame}
\setbeamercolor{background canvas}{bg=} \includepdf{aux/jarosch_f6a.pdf}

\setbeamercolor{background canvas}{bg=} \includepdf{aux/jarosch_f6b.pdf}

\setbeamercolor{background canvas}{bg=} \includepdf{aux/jarosch_f6c.pdf}

\setbeamercolor{background canvas}{bg=} \includepdf{aux/jarosch_f6d.pdf}
\begin{frame}{Decomposition}

\begin{itemize}
\item {\small{}Three sources of wage scarring:\medskip{}
}

\begin{itemize}
\item {\small{}Negotiation rents}{\small \par}
\item {\small{}Position in job ladder (employer itself)}{\small \par}
\item {\small{}Human capital/experience\medskip{}
}{\small \par}
\end{itemize}
\item {\small{}Jarosch uses the model to ``turn on'' each component sequentially}{\small \par}
\end{itemize}
\end{frame}
\setbeamercolor{background canvas}{bg=} \includepdf{aux/jarosch_f7.pdf}\end{document}

